% =====================================================================
% Split Sequence Attribute Tables across SatSOT, SV248S, and OOTB
% Requires: \usepackage{booktabs}, \usepackage{array}, \usepackage{makecell}
% Optional (recommended for two-column layouts): \usepackage{adjustbox}
%
% Revision notes (w.r.t. original tables):
%   1. Occlusion-related attributes (POC/FOC in SatSOT, STO/LTO/CO in
%      SV248S, PO/FO in OOTB) are consolidated into a single unified
%      OCC row in Table 1, because SV248S defines occlusion along a
%      *temporal* axis (short- vs. long-term) whereas SatSOT and OOTB
%      define it along a *spatial* axis (partial vs. full). The two
%      axes are not comparable at the sequence-attribute level, so we
%      evaluate any-occlusion as the only safely shared label.
%   2. ARC and OON are separated into two distinct rows in Table 2
%      (unique attributes) because ARC measures *temporal change* of
%      aspect ratio while OON measures *static extremeness* of the
%      aspect ratio -- they probe different tracker capabilities.
%   3. The unified abbreviation for SV248S BCL is renamed from "INV"
%      to "IBG" (Indistinguishable from Background) to avoid confusion
%      with the SV248S *frame-level* INV flag, of which BCL is the
%      *sequence-level* aggregation.
%   4. ARC definition now reflects the SatSOT specification: the ratio
%      of the current-frame bbox aspect ratio to the first-frame bbox
%      aspect ratio is outside [0.5, 2].
% =====================================================================


% ---------------------------------------------------------------------
% Table 1: Shared attributes (annotated in at least two datasets)
% ---------------------------------------------------------------------
\begin{table*}[t]
\centering
\caption{Shared sequence attributes across SatSOT, SV248S, and OOTB.
Rows list attributes that are annotated in at least two of the three
datasets. The last three columns give the corresponding per-dataset
label (``---'' if not annotated).
Dataset-specific abbreviations appearing in the last three columns are:
\textbf{IPR} = In-Plane Rotation,
\textbf{DS} = Dense Similarity,
\textbf{PO} = Partial Occlusion,
\textbf{FO} = Full Occlusion,
\textbf{STO} = Short-Term Occlusion,
\textbf{LTO} = Long-Term Occlusion,
\textbf{CO} = Continuous Occlusion,
\textbf{SA} = Similar Appearance.}
\label{tab:shared_attributes}
\small
\setlength{\tabcolsep}{4pt}
\renewcommand{\arraystretch}{1.15}
\resizebox{\textwidth}{!}{%
\begin{tabular}{llp{7.4cm}ccc}
\toprule
\textbf{Abbr.} & \textbf{Full Name} & \textbf{Definition} & \textbf{SatSOT} & \textbf{SV248S} & \textbf{OOTB} \\
\midrule
BC   & Background Clutter       & Background has similar appearance (texture/color) to the target                                       & BC  & ---                        & BC  \\
IV   & Illumination Variation   & Illumination around the target changes significantly                                                  & IV  & IV                         & IV  \\
ROT  & Rotation (in-plane)      & Target rotates in the image plane ($\geq 30^{\circ}$ for SV248S)                                      & ROT & IPR                        & IPR \\
OCC  & Occlusion$^{1}$          & Target is occluded by scene structures (any degree, any duration)                                     & POC $\cup$ FOC & STO $\cup$ LTO $\cup$ CO & PO $\cup$ FO  \\
SOB  & Similar Object           & Nearby objects share shape/type/appearance with the target ($\leq 2.5\times$ object size for SV248S)  & SOB & DS$^{2}$                   & SA  \\
DEF  & Deformation              & Non-rigid object deformation                                                                          & DEF & ---                        & DEF \\
\bottomrule
\end{tabular}
}
\vspace{2pt}
\begin{flushleft}\footnotesize
$^{1}$ The three datasets annotate occlusion along incompatible axes:
SatSOT and OOTB split it \emph{spatially} into partial vs. full, while
SV248S splits it \emph{temporally} into short-term ($\leq 50$
consecutive OCC frames), long-term ($>50$ consecutive OCC frames), and
continuous (two or more STO/LTO events in one sequence). We therefore
unify them into a single any-occlusion label and treat spatial-degree
and temporal-length evaluations as dataset-unique (Table~\ref{tab:unique_attributes}).
\\
$^{2}$ SV248S DS requires the distractor to be within $2.5\times$ the
target's bounding-box diagonal, whereas SatSOT SOB and OOTB SA do not
specify a spatial threshold. The mapping is therefore conceptual, not
definitional.
\end{flushleft}
\end{table*}


% ---------------------------------------------------------------------
% Table 2: Dataset-unique attributes (annotated in exactly one dataset,
% or split by incompatible definitions that cannot be safely merged)
% ---------------------------------------------------------------------
\begin{table*}[t]
\centering
\caption{Dataset-unique sequence attributes across SatSOT, SV248S, and OOTB.
Rows list attributes that are annotated in exactly one of the three
datasets, or that are nominally related across datasets but whose
definitions are incompatible enough that merging would be misleading
(POC vs.\ STO, FOC vs.\ LTO, ARC vs.\ OON). The last three columns
give the corresponding per-dataset label (``---'' if not annotated).
The dataset-specific abbreviation appearing in the last three columns is:
\textbf{BCL} = Background Cluster ($\geq 10$ frames carrying the SV248S
frame-level INV flag).}
\label{tab:unique_attributes}
\small
\setlength{\tabcolsep}{4pt}
\renewcommand{\arraystretch}{1.15}
\resizebox{\textwidth}{!}{%
\begin{tabular}{llp{7.4cm}ccc}
\toprule
\textbf{Abbr.} & \textbf{Full Name} & \textbf{Definition} & \textbf{SatSOT} & \textbf{SV248S} & \textbf{OOTB} \\
\midrule
\multicolumn{6}{l}{\emph{Occlusion sub-types (kept separate due to incompatible annotation axes; see Table~\ref{tab:shared_attributes}, footnote 1)}} \\
POC  & Partial Occlusion       & Target is partially (but not fully) occluded in the image plane                                                            & POC & ---          & PO  \\
FOC  & Full Occlusion          & Target is fully occluded in the image plane                                                                                & FOC & ---          & FO  \\
STO  & Short-Term Occlusion    & $\leq 50$ consecutive frames carry the OCC flag                                                                            & --- & STO          & --- \\
LTO  & Long-Term Occlusion     & $>50$ consecutive frames carry the OCC flag                                                                                & --- & LTO          & --- \\
CO   & Continuous Occlusion    & Two or more STO/LTO events occur within a single sequence                                                                  & --- & CO           & --- \\
\midrule
\multicolumn{6}{l}{\emph{Aspect-ratio attributes (kept separate: ARC is a temporal-change property, OON is a static-extremeness property)}} \\
ARC  & Aspect Ratio Change     & Ratio of the current-frame bbox aspect ratio to the first-frame bbox aspect ratio is outside $[0.5, 2]$                    & ARC & ---          & --- \\
OON  & Out-of-Normal           & The bounding-box aspect ratio itself is outside $[0.3, 3]$ in the sequence                                                 & --- & ---          & OON \\
\midrule
\multicolumn{6}{l}{\emph{Other dataset-unique attributes}} \\
LQ   & Low Quality             & Image quality is low; target is hard to distinguish                                                                        & LQ  & ---          & --- \\
TO   & Tiny Object             & At least one GT bbox has fewer than 25 pixels                                                                              & TO  & ---          & --- \\
BJT  & Background Jitter       & Background jitter caused by satellite camera shaking                                                                       & BJT & ---          & --- \\
BCH  & Background Change       & Background has noticeable changes in color or texture along the sequence                                                   & --- & BCH          & --- \\
ND   & Natural Disturbance     & Target appearance affected by smog, sand, or clouds                                                                        & --- & ND           & --- \\
IBG  & Indistinguishable Bg.$^{3}$ & Target disappears without occluder (too similar to surroundings) for $\geq 10$ frames                                  & --- & BCL$^{3}$    & --- \\
SM   & Slow Motion             & Target moves slowly ($< 2.2$ pps in SV248S)                                                                                & --- & SM           & --- \\
LT   & Less Textures           & Target has poor texture information, causing discrimination difficulty                                                     & --- & ---          & LT  \\
MB   & Motion Blur             & Target region is blurred due to object or platform motion                                                                  & --- & ---          & MB  \\
IM   & Isotropic Motion        & Nearby objects move with similar magnitude and direction                                                                   & --- & ---          & IM  \\
AM   & Anisotropic Motion      & Nearby objects move with similar magnitude but opposite direction                                                          & --- & ---          & AM  \\
\bottomrule
\end{tabular}
}
\vspace{2pt}
\begin{flushleft}\footnotesize
$^{3}$ In SV248S, INV is defined at the \emph{frame} level (a target is
invisible without an occluder) and BCL is its \emph{sequence}-level
aggregation. We use the unified label IBG (Indistinguishable from
Background) to avoid conflating the frame flag with the sequence
attribute.
\end{flushleft}
\end{table*}
